\chapter{Shadow Rendering: Moment-Based Deep Shadow Maps}\label{ch:shadow-rendering}

\section{Why deep shadow maps}

A translucent tetrahedral volume cannot use a classic depth-test shadow map:
a light ray is not blocked at a single depth, it is \emph{attenuated
continuously}. What a receiver needs is the transmittance function $T(z)$
along the light ray. \code{rmesh-dsm} stores a compact statistical summary
of the absorption along each light ray --- two $\alpha$-weighted depth
moments --- in a per-light cubemap, and reconstructs a conservative $T(z)$
bound at shading time. The design is documented in \code{MSM.md}; an earlier
4-moment trigonometric variant (Peters--Klein style) was removed in favor of
this cheaper 2-moment scheme.

\section{Representation}

One \code{Rgba16Float} cubemap per light (max 16 lights), each texel:
\[
  \big(\, E[\alpha z],\; E[\alpha z^2],\; 0,\; \alpha_{\text{total}} \,\big),
\]
where $z$ is view-space depth normalized to $[0,1]$ over the light's
$[\text{near}, \text{far}]$ range and
$\alpha_{\text{total}} = 1 - T_{\text{total}}$ is the total absorption along
the ray. At $256^3$ this costs 3\,MiB per light (versus 9\,MiB for the
removed 4-moment variant).

\section{Generating the moments}

\code{generate_dsm_for_lights} (\code{crates/rmesh-dsm/src/lib.rs}) runs
per light, \emph{per cubemap face} ($\times 6$):

\begin{enumerate}
\item \textbf{Build the face view-projection} (\code{build_light_vp}):
  point lights get a $90^\circ$ perspective per face (using face
  direction/up pairs matched to WebGPU cube conventions, plus a clip-space
  Y flip); spot lights a perspective with FOV $= 2\arccos(\cos_{\text{outer}})$;
  directional lights an orthographic box.
\item \textbf{Opaque blocker pre-pass} (\code{DsmPrimitivePipeline},
  \code{dsm_primitive.wgsl}): primitives write $(z, z^2, 0, 1)$ ---
  a full-$\alpha$ absorber concentrated at their depth --- with depth
  writes on, so tets behind opaque geometry are culled by the depth test.
\item \textbf{Project + sort}: the forward \code{project_compute_hw}
  pipeline projects all tets through the face VP and emits depth keys, then
  the radix sort orders them back-to-front \emph{for that face}.
  Each face gets its own project + sort because a single per-light radial
  sort is \emph{not} a valid per-face compositing order off-axis --- this
  correctness fix (commit \code{9207ea3}) costs $6\times$ the sort work per
  light but eliminates wrong shadows.
\item \textbf{Interval generation + moment render}
  (\code{DsmPipeline}): the interval renderer's geometry pass runs from the
  light's viewpoint, and \code{dsm_moment_fragment.wgsl} converts each
  tet's depth interval into moment contributions, blended back-to-front
  with premultiplied alpha
  ($\text{src}\cdot 1 + \text{dst}\cdot(1-\alpha_{\text{src}})$) into a
  shared 2D staging texture, which is then copied into the cube face layer.
\end{enumerate}

\paragraph{Moment fragment math.} For a tet crossing the light ray over
$[z_a, z_b]$ with optical depth $\tau$: $\alpha = 1 - e^{-\tau}$, and using
the same termination-distribution weights as the color integral
($\varphi = (1-e^{-\tau})/\tau$, $w_0 = \varphi - e^{-\tau}$,
$w_1 = 1 - \varphi$):
\[
  \text{out} = \big( w_0 z_b + w_1 z_a,\;\; w_0 z_b^2 + w_1 z_a^2,\;\; 0,\;\;
  \alpha \big) \;=\; \big( \alpha E[z],\; \alpha E[z^2],\; 0,\; \alpha \big).
\]

Regeneration is triggered when lights move or appear (e.g.\ a flying flare
projectile), or on demand from the UI.

\section{Atlas and metadata}

\code{DsmAtlas} owns the per-light cubemaps, the staging textures, and a
storage buffer of \code{ShadowLightMeta}: per light, the six per-face
view-projection \emph{W rows only} (enough to compute the receiver's
light-space depth; storing full matrices cost $4\times$ the bytes and forced
materializing all six for dynamic indexing), plus near/far and the light
type.

\section{Sampling at shading time}

The deferred pass (Chapter~\ref{ch:deferred-shading}) calls
\code{evaluate_transmittance(world_pos, light)} per light:

\begin{enumerate}
\item Select the cube face from the major axis of
  $d = p_{\text{world}} - p_{\text{light}}$ and sample the moments with a
  filtering cube sampler.
\item Receiver depth: $w_{\text{clip}} = \text{VP}_{w\,\text{row}} \cdot
  (p_{\text{world}}, 1)$, normalized
  $z = \frac{w_{\text{clip}} - n}{f - n} - \text{bias}$.
\item Un-premultiply: $\mu = \frac{E[\alpha z]}{\alpha}$,
  $\overline{z^2} = \frac{E[\alpha z^2]}{\alpha}$.
\item If $z \le \mu$, the receiver is in front of the mean absorber:
  $T = 1$. Otherwise apply the \textbf{Cantelli (one-sided Chebyshev)
  bound} --- the same inequality behind Variance Shadow Maps:
  \[
    \sigma^2 = \max\!\big(\overline{z^2} - \mu^2,\;
    \sigma^2_{\min}\big), \qquad
    T(z) \;\le\; p_{\max} = \frac{\sigma^2}{\sigma^2 + (z - \mu)^2}.
  \]
\item Reduce light bleeding by remapping
  $\text{linstep}(0.15, 1, p_{\max})$, and floor the result at
  $1 - \alpha_{\text{total}}$ (light that provably gets through).
\end{enumerate}

Constants: depth bias $1.5 \times 10^{-3}$, minimum variance
$4 \times 10^{-4}$, bleed-reduction threshold $0.15$. The bound is exact
only for a two-spike absorber; for smoother distributions it
\emph{over}-shadows, which errs on the visually safe side.

\section{Debugging and validation}

\code{DsmResolvePipeline} (\code{dsm_resolve.wgsl}) reconstructs $T(z)$ from
a moment texture at a UI-chosen query depth --- it must use identical
constants to the deferred shader --- and
\code{record_dsm_resolve_cubemap_cross} lays the six faces out as a
$3\times2$ cross for inspection. The deferred pass also has dedicated debug
modes for average transmittance, cube-face selection, and raw DSM moments.
A CPU-reference comparison lives in
\code{crates/rmesh-dsm/tests/probe_vs_cpu_test.rs}.

References: Donnelly \& Lauritzen, \emph{Variance Shadow Maps} (2006);
Peters \& Klein, \emph{Moment Shadow Mapping} (2015); Lokovic \& Veach,
\emph{Deep Shadow Maps} (2000).

